\section{Generic Reverse-AD of Reduce with Invertible Operator}
My code validates and is provided below:
\begin{lstlisting}[language={futhark}]
def vjp2_red_inv 't 't_lft [n]
               (op : t -> t -> t)
               (op_lft : t_lft -> t_lft -> t_lft)
               (ne: t_lft)
               (cfwd : t -> t_lft)
               (cbwd : t_lft -> t)
               (op_inv : t_lft -> t -> t)
               (as : [n]t)
               (rs_bar: t)
             : (t, [n]t) =
  -- please replace the dummy code below with your correct implementation
  let t_prim = map cfwd as |> reduce_comm (op_lft) ne
  let y = cbwd t_prim
  let as_bar = map (\a ->
    let p' = op_inv t_prim a
    in vjp (\x -> op x  p') a rs_bar
  ) as
  in (y, as_bar)
\end{lstlisting}

If we assume no special cases, then the differentiation of an arbitrary reduce
operation would lead to two scans and a map call. We are able to circumvent
this, such that we instead have a single map call, and then whatever is the
differentiation of the operator, which often is just constant time. This means,
that we are able to go from work O($n$) and span O($\lg n$) to work O($n$) and
span O($1$) and we should also have reduced the constant factor by now only
doing a single map call.
\begin{table}[H]
\centering
\caption{Benchmark runtimes and speedups for Mul and Sop groups.}
\label{tab:generic}
\begin{tabular}{l l cc cc}
\toprule
Operator & Method & \multicolumn{2}{c}{Runtime (\textmu s)} & \multicolumn{2}{c}{Overhead from primal} \\
\cmidrule(lr){3-4} \cmidrule(lr){5-6}
      &        & 50M elements & 500M elements & 50M elements & 500M elements \\
\midrule
\multirow{4}{*}{Mul} 
 & primal & 296   & 2734   & -    & - \\
 & ppad   & 3414  & 33662  & 11.53    & 12.31 \\
 & vjp2   & 901   & 8603   & 3.04    & 3.15 \\
 & manual & 920   & 8717   & 3.11 & 3.19 \\
\midrule
\multirow{4}{*}{Sop} 
 & primal & 594   & -      & -    & - \\
 & ppad   & 6214  & -      & 10.46   & - \\
 & vjp2   & 3864  & -      & 6.51   & - \\
 & manual & 1450  & -      & 2.44 & - \\
\bottomrule
\end{tabular}
\end{table}

The runtimes can be seen in Table~\ref{tab:generic}, which shows that in the
Mul case the runtimes of vjp2 and manual are almost equivalent. However, in the
Sop case the manual implementation is significantly faster. This is most likely
because futhark has this optimisation built in for the built in operators such
as multiplication, but not for the sop case.
